\justifying
\NSFCBodyText

\subsubsection{申请人正在承担的与本项目相关的科研项目}

\subsubsubsection{多模态融合技术研究及其在新疆产业集群中的应用}

资助机构与项目类别：自治区天池英才引进计划（青年博士人才）。该项目无独立项目批准号，依据为项目任务书。获资助金额 50 万元，起止年月为 2025 年 2 月至 2028 年 2 月，申请人为主持人，负责总体研究方案、关键模型设计与成果组织。

与本项目的关系及负责内容：该项目面向多模态融合技术及其在区域产业中的应用，研究范围较宽；本申请聚焦于遥感图文检索这一具体任务中"描述不完备条件下地理实体关系证据的状态判定"这一科学问题，研究对象、拟解决的关键问题、使用数据与评测方式均与该项目不同。二者的具体区分结论需以任务书逐项比对为准，当前核查状态为\textbf{【待补 F-MAP-03：本申请与天池英才项目在研究对象、科学问题、数据、方法、经费与成果六个维度的重复性核查结论】}。

\subsubsubsection{数据驱动下基于互补型生成对抗网络与神经统计模型的公路沙害智能预警研究}

资助机构与项目类别：国家自然科学基金地区科学基金项目，批准号 62466056。项目获资助总额 32 万元，基金委系统登记单位为新疆大学，起止年月为 2025 年 1 月至 2028 年 12 月。申请人为参与者，负责数据的收集整理。

与本项目的关系及负责内容：该项目面向公路沙害智能预警，与本申请的遥感图文检索任务在科学问题与技术路线上均不相同；申请人在该项目中承担数据收集整理工作，与本申请的职责不重叠。

\subsubsubsection{可学习乘积量化与状态空间模型协同的长上下文高效注意力机制研究}

资助机构与项目类别：CAAI---蚂蚁科研基金，批准号 202604140006。项目获资助总额 15 万元，起止年月为 2026 年 2 月至 2027 年 2 月。申请人为参与者，负责高效注意力机制研究。

与本项目的关系及负责内容：该项目面向长上下文高效注意力机制，属模型效率方向的基础方法研究，不涉及遥感图文检索、地理实体关系建模与检索评测协议；其执行期在本申请执行期开始前后结束。

\subsubsubsection{正在承担项目清单的核定说明}

上述列示依据申请人已确认的项目信息编制。由于"正在承担"是申报时点的状态声明，此前轮次确认的列示口径中部分项目将于本申请执行期开始前结束，不能直接沿用。本申请提交时，正在承担项目的最终清单为\textbf{【待补 F-MAP-08：2027 年申报时点经基金委系统核对后的正在承担项目最终清单，含项目类别、批准号、名称、资助金额、起止年月与本人负责内容】}，届时将按系统记录逐项核对后据实填写。

\subsubsection{主要参与者正在承担的与本项目相关的科研项目}

本项目主要参与者的名单与分工为\textbf{【待补 F-MAP-01：本项目主要参与者姓名、职务职称、在本项目中的具体分工】}。申请人已确认，拟参与的研究生均\textbf{无正在承担且与本申请相关的科研项目}。

此处需要交代一点：共用申请人信息中登记的研究生分工安排形成于另一研究方向的申报，与本申请的遥感图文检索方向不匹配，因此本申请不沿用该分工描述，而是将参与者的具体分工列为待核定事项，待申请人按本项目研究内容重新确认后填写。

\subsubsection{与本项目的区别与联系}

申请人正在承担的项目中，与本申请在技术领域上存在一定关联的是多模态融合方向的天池英才项目，其余项目在研究对象与科学问题上与本申请没有实质交叉。本申请与前者共享多模态表征学习这一技术基础，但研究对象从"区域产业应用中的多模态融合"收敛到"遥感图文检索中描述不完备条件下的关系证据状态判定"，科学问题、验证数据与成果形式均独立设置。

申请人承诺：不将立项前已形成的论文、模型或代码认定为本项目的阶段成果，不将同一成果重复列为多个项目的产出；成果归属边界以各项目任务书与成果标注为准。重复性核查的最终结论以 F-MAP-03 与 F-MAP-08 的核定结果为准。
